% Đầu sách: đọc tài liệu này thế nào. Không đánh số, để Phần I bắt đầu từ Chương 1.

\chapter*{Đọc tài liệu này thế nào}
\addcontentsline{toc}{chapter}{Đọc tài liệu này thế nào}

Ai cũng nên đọc mục \emph{Làm quen với Kriky}. Sau đó, thầy cô và các em chỉ cần đọc phần của mình.

\begin{bangtra}{L{4.2cm}L{3.2cm}Y}{\dau{Bạn là} & \dau{Đọc} & \dau{Trong đó}}
  Giáo viên & Phần I, trang~\pageref{phan:gv} & Lập lớp, soạn đề cùng Kriky, phát hành, xem
  kết quả \\
  Học sinh & Phần II, trang~\pageref{phan:hs} & Đăng nhập, làm bài, hỏi Kriky câu sai, làm
  lại \\
  Ai cũng nên đọc & Làm quen, trang~\pageref{chuong:lam-quen} & Kriky là ai, cần chuẩn bị gì \\
\end{bangtra}

\section{Các ký hiệu trong sách}

\begin{bangtra}{L{4.6cm}Y}{\dau{Thấy trong sách} & \dau{Nghĩa là}}
  \nut{Duyệt đề} & Một nút trên màn hình. Bấm vào đó. \\
  \nhan{Tên lớp} & Một ô, một nhãn hay một dòng chữ trên màn hình. \\
  \go{12A} & Chữ gõ vào, đúng từng chữ. \\
  \sobuoc{1} & Một bước. Số trên ảnh khớp với số ở bước. \\
  \leavevmode{\color{correct}\faCheckCircle\enspace\vua Xong khi:} & Thấy điều này là việc đã xong. \\
  \leavevmode{\color{krikyorangedeep}\vua Viền cam trên ảnh} & Vùng cần để ý trên màn hình. \\
  Ô xanh nhạt lệch phải & Một câu mẫu để nhắn cho Kriky. \\
\end{bangtra}

\machnho{Ô vàng có mặt Kriky là mẹo để làm nhanh hơn.}

\luuy{Ô nâu là điều cần để ý, bỏ qua dễ mất công làm lại.}

\neuloi{Ô đỏ chỉ cách xử lý khi màn hình báo lỗi.}
